% arxiv-template.tex -- pandoc template for the arXiv/preprint version.
%
% Consumed by: quarto render stpois_sj.qmd --to latex \
%   --template ../arxiv/arxiv-template.tex --output stpois_arxiv.tex
%
% Defines Stata Journal macro compatibility so the body produced from
% stpois_sj.qmd compiles without the SJ style files.

\documentclass[12pt,a4paper]{article}

\usepackage[T1]{fontenc}
\usepackage[utf8]{inputenc}
\usepackage{amsmath}
\usepackage{amssymb}
\usepackage{bm}
\usepackage[round,authoryear]{natbib}
\usepackage[colorlinks=true,linkcolor=blue,citecolor=blue,urlcolor=blue]{hyperref}
\usepackage{microtype}
\usepackage[margin=1in]{geometry}
\usepackage{booktabs}
\usepackage{xcolor}

% ── Stata Journal macro layer ─────────────────────────────────────────────
% Load StataCorp's official stata.sty (same file the SJ build uses) so the
% Stata log tables, box-drawing rules (\HLI/\TOPT/\PLUS/\BOTT/\VBAR),
% command-syntax boxes (stsyntax), and helper macros (\stcmd, \stref,
% \varlist, \optional, \num, \oom, \hangpara, ...) render identically to the
% Stata Journal PDF instead of being stubbed out.
\usepackage{stata}

% Article tag (SJ-only macro, used in the keywords block; not in stata.sty)
\def\inserttag{stpois}

% Pandoc compact-list boilerplate
\providecommand{\tightlist}{%
  \setlength{\itemsep}{0pt}\setlength{\parskip}{0pt}}

% ── Document metadata ─────────────────────────────────────────────────────

\title{stpois: Fast Poisson event-history regression with
high-dimensional fixed effects}

\author{%
  Andreas Ljungström\\
  \small Swedish Institute for Social Research (SOFI), Stockholm University\\
  \small\texttt{andreas.ljungstrom@sofi.su.se}%
}

\date{}

% ── Document body ─────────────────────────────────────────────────────────

\begin{document}

\maketitle

\begin{abstract}
\noindent Poisson regression on episode-split survival data is a
workhorse of event-history analysis, but two practical obstacles limit
its use with full-population register data: models with rich
fixed-effect structures become slow or infeasible to estimate with
explicit dummy variables, and continuous covariates prevent the
collapsing of individual records to risk cells that otherwise makes such
models fast. This article introduces the community-contributed command
\stcmd{stpois}, which fits Poisson event-history models directly on
\stcmd{stset} data through three estimation paths, all of which return
maximum-likelihood estimates that agree with \stcmd{poisson} to at least
\(10^{-6}\): a standard path that wraps \stcmd{poisson} with automatic
exposure; a high-dimensional fixed-effects path implemented via
iteratively reweighted alternating projections in Mata; and a fast path
that accelerates estimation through the cell structure of the
categorical covariates. The fast path computes each Newton iteration
from exponentially tilted within-cell moments in a single pass over the
microdata. Coefficients, log likelihoods, and observed-information,
robust, and cluster--robust standard errors agree with the corresponding
\stcmd{poisson} model on every path to at least \(10^{-6}\) (typically
\(\sim 10^{-8}\)); factor-variable interactions and weights are
supported throughout, the two acceleration mechanisms can be combined,
and postestimation predictions after absorption include the recovered
fixed effects. Benchmarks against \stcmd{streg}, \stcmd{poisson}, and
\stcmd{ppmlhdfe} document the speed gains in three settings:
episode-split workflows, all-categorical rate models, and models with
high-dimensional fixed effects.
\end{abstract}

\medskip\noindent\textbf{Keywords:} \inserttag, stpois, event-history
analysis, Poisson regression, piecewise exponential model,
high-dimensional fixed effects, register data

\bigskip

\section{Introduction}

\label{sec:intro}

The equivalence between piecewise-constant hazard models and Poisson
regression on episode-split data is one of the older and more useful
results in event-history analysis
\citep{holford1980, laird1981, whitehead1980}. Splitting each subject's
time at risk into episodes, treating the event indicator in each episode
as a Poisson count, and entering the log of exposure time as an offset
yields maximum-likelihood estimates of a proportional-hazards model with
a hazard that is constant within episodes. In Stata, this model can be
fitted with \stcmd{poisson} after manually generating the exposure
variable, or---for the single-episode exponential case---with
\stcmd{streg, distribution(exponential)}. The approach is standard in
epidemiology \citep{clayton1993, carstensen2007} and is covered at
length by \citet{rabe2022}.

Two practical obstacles arise when this framework meets modern
administrative data. Full-population registers of the kind available in
the Scandinavian countries routinely produce analysis files with
millions or tens of millions of episode records, particularly after
splitting on calendar time, age, and duration. First, models with rich
fixed-effect structures---birth cohorts, municipalities, detailed
occupational codes, or their interactions---become slow or infeasible
when the fixed effects are entered as explicit dummy variables, since
estimation time and memory grow quadratically with the parameter count.
The \stcmd{ppmlhdfe} command \citep{correia2020} solves this problem for
general Poisson models, building on the alternating-projections approach
of \citet{guimaraes2010}; it depends on the \stcmd{ftools} and
\stcmd{reghdfe} packages, so deploying it means installing that chain,
which can add friction in the locked-down secure environments where
register data must typically be analyzed. Second, and less widely
appreciated, the Poisson likelihood for purely categorical covariates
depends on the data only through the sums of events and exposure within
covariate cells. Estimation on the collapsed cell data is then
\emph{exact} and can reduce the per-iteration operation count by orders
of magnitude when the number of cells is far smaller than the number of
records. A single continuous covariate breaks this property, and
analysts face an unattractive choice between discretizing the covariate
and forgoing the speed advantage entirely.

The command described in this article, \stcmd{stpois}, addresses both
obstacles within a single interface for \stcmd{stset} data. It offers
three estimation paths, all of which return maximum-likelihood estimates
and standard errors that agree with \stcmd{poisson} to the tolerances
documented in section~\ref{sec:validation}:

\begin{enumerate}
\item A \textbf{standard path} that wraps \stcmd{poisson}, using the event indicator \stcmd{\_d} as outcome and time at risk \stcmd{\_t} $-$ \stcmd{\_t0} as exposure.  Coefficients and standard errors are identical to \stcmd{streg, distribution(exponential)}.
\item A \textbf{high-dimensional fixed-effects path}, \stcmd{absorb()}, implementing iteratively reweighted least squares with alternating-projections demeaning in Mata---the same mathematical approach as \stcmd{ppmlhdfe}, but self-contained with no dependencies.  Absorbed terms may be single variables or interactions like \stcmd{region\#period}.  Observed-information, robust, and cluster--robust standard errors agree with those from the corresponding dummy-variable model to at least $10^{-6}$.
\item A \textbf{fast path}, option \stcmd{fast}, which accelerates exact estimation through the cell structure of the categorical covariates.  Newton iterations on the full likelihood are computed from exponentially tilted within-cell moments of the individual-level covariates, accumulated in a single pass over the microdata per iteration; all remaining algebra runs on the cells.  When every covariate is categorical, the iterations run entirely on the collapsed cell sums.  The fast path requires at least one categorical term to define the cells.
\end{enumerate}

Factor-variable interactions---continuous\#continuous,
continuous\#categorical, and categorical\#categorical---are supported on
all paths, as are frequency, importance, and sampling weights, with
weighted results matching the weighted \stcmd{poisson} model to the same
tolerance. The \stcmd{fast} and \stcmd{absorb()} paths may be combined:
the absorbed fixed effects are then concentrated out of the
cell-accelerated likelihood. Section~\ref{sec:model}
sets out the modeling framework, section~\ref{sec:fast}
describes the fast path, and section~\ref{sec:hdfe}
describes the fixed-effects algorithm.
Sections~\ref{sec:command}
and~\ref{sec:examples} document the syntax and
illustrate the command, and section~\ref{sec:validation}
verifies the numerical claims and benchmarks the command against
\stcmd{streg} and \stcmd{poisson}.

\section{Poisson event-history regression}

\label{sec:model}

Consider \stcmd{stset} data in which record \(i\) contributes an event
indicator \(d_i \in \{0, 1\}\) (Stata's \stcmd{\_d}) and a time at risk
\(t_i = \mbox{\stcmd{\_t}} - \mbox{\stcmd{\_t0}}\). Under a
proportional-hazards model with hazard rate
\(\lambda_i = \exp(\mathbf{x}_i'\bm{\beta})\) constant within records,
the log likelihood for right-censored data is, up to a constant,

\begin{equation*}
\ell(\bm{\beta}) = \sum_i \left\{ d_i \, \mathbf{x}_i'\bm{\beta}
  + d_i \log t_i - t_i \exp(\mathbf{x}_i'\bm{\beta}) \right\},
\end{equation*}

in which the \(d_i \log t_i\) term is constant in \(\bm{\beta}\); this
is proportional to the likelihood of a Poisson model for \(d_i\) with
mean \(\mu_i = t_i \exp(\mathbf{x}_i'\bm{\beta})\) and offset
\(\log t_i\). Hence any software that fits Poisson regression with an
offset fits this hazard model, and duration dependence or period effects
are accommodated by episode splitting (\stcmd{stsplit}) followed by
indicator variables for the resulting time bands
\citep{holford1980, laird1981}. The Poisson ``counts'\,' here are 0/1
event indicators; nothing requires the outcome to be a genuine count,
and the pseudo-maximum-likelihood results of \citet{gourieroux1984}
imply that the coefficient estimates are consistent whenever the
log-linear mean is correctly specified, with robust standard errors
available in the usual way \citep[see also][]{santos2006}.

\stcmd{stpois} fits exactly this model on the current \stcmd{stset}
data. On the standard path, it constructs the exposure variable, calls
\stcmd{poisson}, and posts results with survival-flavored headers
(numbers of subjects and failures, time at risk). The point of the
command, however, is not the convenience wrapper but the two departures
from it: the fast path (section~\ref{sec:fast}) and the
fixed-effects path (section~\ref{sec:hdfe}).

\subsection{Sufficiency of cell sums under categorical
covariates}\label{sufficiency-of-cell-sums-under-categorical-covariates}

\label{sec:suff}

When all covariates are categorical, the covariate vector takes a finite
number of distinct values, indexing cells \(j = 1, \ldots, J\). The log
likelihood above then depends on the data only through the within-cell
sums \(D_j = \sum_{i \in j} d_i\) and \(T_j = \sum_{i \in j} t_i\):

\begin{equation*}
\ell(\bm{\beta}) = \sum_j \left\{ D_j\, \mathbf{x}_j'\bm{\beta}
  - T_j \exp(\mathbf{x}_j'\bm{\beta}) \right\} + \text{const},
\end{equation*}

\((D_j, T_j)\) are sufficient statistics, and estimation on the \(J\)
collapsed cells reproduces the individual-level maximum-likelihood
estimates exactly---coefficients, standard errors, and all. With
millions of episode records but a few hundred cells (period \(\times\)
education \(\times\) region, say), the gain is dramatic, and this is the
standard trick behind register-based rate modeling
\citep{clayton1993, dickman2004}. A continuous covariate destroys the
sufficiency because \(\exp(\gamma x_i)\) no longer aggregates:
\(\sum_{i \in j} t_i \exp(\gamma x_i) \neq T_j \exp(\gamma \bar{x}_j)\)
by strict convexity of the exponential, where
\(\bar{x}_j = T_j^{-1}\sum_{i \in j} t_i x_i\) is the exposure-weighted
within-cell mean (the two sides coincide only in the degenerate case of
no within-cell variation in \(x\)). The fast path generalizes the
collapsing idea so that it delivers the exact maximum-likelihood
estimator in the presence of continuous covariates as well.

\section{Fast exact estimation: The fast option}

\label{sec:fast}

\subsection{Term classification and
cells}\label{term-classification-and-cells}

With \stcmd{fast}, \stcmd{stpois} classifies each expanded term of the
varlist by its pieces. Terms built entirely from factor-variable
pieces---\stcmd{i.edu}, \stcmd{i.edu\#i.region}---are constant within
cells and enter the model at the cell level, with coefficient vector
\(\bm{\delta}\) on the cell-level design \(\mathbf{w}_j\). Terms
involving at least one continuous piece---\stcmd{age},
\stcmd{c.age\#c.income}, \stcmd{c.age\#i.edu}---enter at the individual
level, with coefficient vector \(\bm{\gamma}\) on the individual-level
design matrix \(\mathbf{X}_i\) (the columns of the expanded model matrix
corresponding to terms with at least one continuous piece), so that
\(\mu_i = t_i \exp(\bm{\gamma}'\mathbf{X}_i + \bm{\delta}'\mathbf{w}_{j(i)})\),
where \(j(i)\) denotes the cell of record \(i\). The cells \(j\) are the
cross-classification of the categorical variables appearing in the
cell-level terms; at least one cell-level term is required.

\subsection{Newton iterations on tilted cell
moments}\label{newton-iterations-on-tilted-cell-moments}

The score and observed information of the Poisson likelihood decompose
into fixed cell sums and \emph{exponentially tilted within-cell moments}
\citep{jensen1995} of the individual-level covariates:

\begin{equation*}
M_j = \sum_{i \in j} \mu_i, \qquad
\mathbf{m}_j = \sum_{i \in j} \mu_i \mathbf{X}_i, \qquad
\mathbf{Q} = \sum_i \mu_i \mathbf{X}_i \mathbf{X}_i'.
\end{equation*}

The score is
\(\left(\sum_i d_i\mathbf{X}_i - \sum_j \mathbf{m}_j,\; \sum_j \mathbf{w}_j (D_j - M_j)\right)\)
and the information matrix has blocks \(\mathbf{Q}\),
\(\sum_j \mathbf{m}_j \mathbf{w}_j'\), and
\(\sum_j M_j \mathbf{w}_j \mathbf{w}_j'\). The \(\mu_i\)-weighting in
these quantities is the standard Poisson information weight:
\citet[sec.~2.5]{mccullagh1989} derive that the IRLS working weight for
observation \(i\) is \(W_i = [V(\mu_i)(d\eta_i/d\mu_i)^2]^{-1}\); with
variance function \(V(\mu)=\mu\) and log-link derivative
\(d\eta/d\mu = \mu^{-1}\) this gives \(W_i = \mu_i\), so \(M_j\),
\(\mathbf{m}_j\), and \(\mathbf{Q}\) are \(W\)-weighted sums of the
exposure, the individual-level covariates, and their outer products.
Because the log link is canonical for Poisson, the expected and observed
information coincide \citep[p.\,53]{mccullagh1989}, so Newton--Raphson
and Fisher scoring reduce to the same algorithm and the step direction
is algebraically exact. Each Newton iteration therefore requires one
pass over the microdata to accumulate the tilted moments---at cost
\(O(np^2)\), with \(p\) the number of \emph{individual-level} columns
only---after which everything happens on the \(J\) cells. A full-model
Newton step in \stcmd{poisson}, by contrast, costs \(O(n(p+k)^2)\) with
\(k\) the number of cell-level parameters, so the per-iteration
floating-point saving is a factor of roughly \((p+k)^2/p^2\), and it
grows with the richness of the categorical structure. Step-halving on
the log likelihood guards the updates, and, writing
\(\bm{\theta} = (\bm{\gamma}', \bm{\delta}')'\) for the stacked
parameter vector, convergence is declared when
\(\max_\ell |\Delta\theta_\ell| / (1 + |\theta_\ell|)\), evaluated at
the updated parameter vector, falls below \stcmd{tolerance()}.

When \emph{every} term is cell-level (\(p = 0\)), the sufficiency result
of section~\ref{sec:suff} applies directly and the
estimator depends on the data only through the collapsed sums
\((D_j, T_j)\). In this case, with no absorbed fixed effects and the
default observed-information VCE, \stcmd{stpois} routes estimation to
Stata's compiled \stcmd{collapse} followed by \stcmd{poisson} on the
cell sums: the fit is exact by sufficiency and faster than iterating in
Mata over the microdata. The robust and cluster--robust variance
estimators need the per-observation scores---the squared cell residual
is not the cell sum of squared residuals---so under \stcmd{vce(robust)}
or \stcmd{vce(cluster)} the pure cell-level model instead runs a Newton
iteration on \((D_j, T_j)\) at \(O(Jk^2)\) per step, independent of the
number of episode records.

This is Newton on the exact likelihood: the score/information
decomposition involves no linearization or discretization, so the
estimator is algebraically the full-model maximum-likelihood estimator,
and the two maximizers differ only at their iterative convergence
tolerance. In practice, coefficients, the log likelihood, and the
observed-information standard errors agree with the full \stcmd{poisson}
model to at least \(10^{-6}\) (typically \(\sim 10^{-8}\)), and robust
and cluster--robust variance estimators are computed from the same
per-observation scores as their \stcmd{poisson} counterparts, with the
HC1 and \(G/(G-1)\) small-sample corrections (\(G\) the number of
clusters). Because the continuous coefficients are estimated from the
full individual-level likelihood, they are identified---with full
efficiency---even when the continuous covariates do not vary across
cells at all.

\subsection{Implementation}\label{implementation}

The data handling is designed for large datasets, adopting the strategy
popularized by the \stcmd{ftools} project \citep{correia2016ftools}:
group structures are computed from a single \stcmd{order()} sort in
Mata, segment boundaries are found by comparing adjacent sorted rows,
and all within-group sums---the cell aggregates, the tilted moments, and
the within-cluster score sums of the cluster--robust variance
estimator---are running-sum differences over the sorted data, each an
\(O(n)\) pass independent of the number of groups. Interaction columns
are constructed directly from the raw variables in Mata. In this Mata
engine---used whenever the model has continuous covariates, and for the
robust and cluster--robust cell-level fits---there are no
\stcmd{preserve}/\stcmd{collapse} round-trips and no temporary variables
in the dataset; beyond the single sort permutation, memory overhead is
one copy of the loaded columns. (The one exception is the purely
categorical model under the default observed-information VCE, which is
delegated to Stata's \stcmd{collapse} and \stcmd{poisson} as described
in section~\ref{sec:fast}.) With weights, every
likelihood contribution, cell sum, and tilted moment is
weight-multiplied, so the weighted estimates remain the exact weighted
maximum-likelihood estimates.

\subsection{Combining fast with absorbed fixed
effects}\label{combining-fast-with-absorbed-fixed-effects}

\label{sec:combo}

\stcmd{fast} may be combined with \stcmd{absorb()}
(section~\ref{sec:hdfe}). The absorbed effects are then
concentrated out of the likelihood by closed-form updates: for Poisson,
the profile-maximizing effect of absorbed group \(g\) given the
remaining parameters is
\(\alpha_g = \ln(\sum_{i \in g} d_i / \sum_{i \in g} \mu_i)\), and
Gauss--Seidel cycles of this update across the absorbed terms are
observed in practice to converge to the profile optimum; formal
convergence guarantees for the nonlinear profile Gauss--Seidel are not
available---the approach is used by the R package \stcmd{FENmlm}
\citep{berge2018} and its successor \stcmd{fixest}. Each Newton step is
then taken on the profile likelihood: the design columns
(individual-level and cell-level) are weighted-demeaned within the
absorbed groups, and the step uses the demeaned score and information.
Given that the inner fixed-effect solve has converged to the profile
optimum, this is, by the Frisch--Waugh--Lovell logic applied through the
IRLS local linearization, the profile Newton direction and, at
convergence, the partitioned variance matrix. Step-halving re-solves the
fixed effects at every trial point, so the profile likelihood increases
monotonically. Coefficients and all three variance estimators agree with
\stcmd{poisson} carrying explicit dummies for both the cell-level terms
and the absorbed effects to at least \(10^{-6}\).

\section{High-dimensional fixed effects: absorb()}

\label{sec:hdfe}

The \stcmd{absorb(}{\it terms\/}\stcmd{)} option fits the exact model
with one or more sets of fixed effects without creating dummy variables,
using iteratively reweighted least squares
\citep[sec.~2.5]{mccullagh1989} (IRLS) combined with weighted
alternating-projections demeaning---the approach introduced to Stata by
\citet{guimaraes2010} and refined by \citet{correia2020}. Each absorbed
term is a variable or an interaction such as \stcmd{region\#period},
absorbed as the cross-classification of its variables. Each IRLS
iteration:

\begin{enumerate}
\item computes the working variable $z_i = \eta_i + (d_i - \mu_i)/\mu_i$ and weights $\omega_i = \mu_i$, where $\eta_i = \mathbf{x}_i'\bm{\beta}$ is the current linear predictor excluding the offset and $\mu_i = t_i\exp(\eta_i)$;
\item demeans $z$ and each covariate within the groups of each absorbed term in turn, iterating this Gauss--Seidel cycle to convergence (weighted alternating projections), which is guaranteed to converge for the linear demeaning problem \citep{halperin1962, gaure2013} though at a rate that deteriorates as the absorbed dimensions become more correlated;
\item regresses the demeaned $\tilde{z}$ on the demeaned $\tilde{\mathbf{X}}$ by weighted least squares to update $\bm{\beta}$; and
\item recovers the fixed-effect contribution as the complement of the projection and updates $\eta$.
\end{enumerate}

Convergence is declared when
\(\max_\ell |\Delta\beta_\ell| / (1 + |\beta_\ell|)\), evaluated at the
updated coefficient vector, falls below \stcmd{tolerance()} (default
\(10^{-8}\)). Sort permutations and group boundaries for each absorbed
term are precomputed once, so every demeaning pass is a single \(O(n)\)
sweep independent of the number of levels. The implementation is pure
Mata with no dependencies, and the non-absorbed varlist may include
factor variables and interactions.

\begin{sloppypar}
Standard errors come from the partitioned information at convergence:
$\mathbf{V} = (\tilde{\mathbf{X}}'\bm{\Omega}\tilde{\mathbf{X}})^{-1}$,
with $\bm{\Omega} = \mathrm{diag}(\omega_i)$, for
\stcmd{vce(oim)}, the HC1 sandwich with meat
$\sum_i \mathbf{s}_i \mathbf{s}_i' \cdot n/(n-1)$ and scores
$\mathbf{s}_i = (d_i - \mu_i)\tilde{\mathbf{x}}_i$ for
\stcmd{vce(robust)}, and
the $G/(G-1)$-corrected clustered sandwich with scores summed within
clusters for \stcmd{vce(cluster)}.  By the
Frisch--Waugh--Lovell logic applied through the IRLS local linearization, these are not linear approximations but the
partitioned estimators; given that the fixed effects are solved to
tolerance, all three match the corresponding dummy-variable
\stcmd{poisson} output to within $10^{-6}$ (typically $\sim 10^{-8}$),
which the test suite verifies (section~\ref{sec:validation}).  The model test reported after
\stcmd{absorb()} is a Wald test of the non-absorbed coefficients,
$\hat{\bm{\beta}}'\mathbf{V}^{-1}\hat{\bm{\beta}}$, which matches the
joint Wald test on the dummy-variable model; a naive likelihood-ratio
test formed from the reported profile log likelihood against a null
without the fixed effects would not be valid, since the test's degrees of
freedom would ignore the absorbed effects.
\end{sloppypar}

Groups of an absorbed term in which no events occur are separated in the
maximum-likelihood sense---their fixed effects diverge to
\(-\infty\)---and \stcmd{stpois} drops such observations before
estimation with a message, mirroring the practice of \stcmd{ppmlhdfe}
\citep[which treats separation in far greater generality]{correia2020}.
The header statistics (subjects, failures, time at risk) are computed on
the final estimation sample after any such drops, so they always agree
with \stcmd{e(N)} and \stcmd{e(sample)}.

\begin{sloppypar}
At convergence, the total fixed-effect contribution
$\hat{\eta}_i - \mathbf{x}_i'\hat{\bm{\beta}}$ is stored in a
variable---\stcmd{\_stpois\_fe} by default, or a name supplied in the
\stcmd{d()} option, mirroring \stcmd{ppmlhdfe}'s \stcmd{d()} and the
\stcmd{fixef()} accessor of \stcmd{fixest}.  \stcmd{predict} adds this
variable to the linear predictor, so the full set of postestimation
statistics---linear predictor, expected events, hazard, and
survival---is available after \stcmd{absorb()} and reproduces the
corresponding dummy-variable \stcmd{poisson} predictions.
\end{sloppypar}

\section{The stpois command}

\label{sec:command}

\subsection{Syntax}\label{syntax}

\begin{stsyntax}
stpois
    \varlist\
    \optif\
    \optin\
    \optweight\
    \optional{,
    irr
    nolog
    \underbar{l}evel(\num)
    \underbar{nocons}tant
    vce({\it vcetype\/})
    \underbar{r}obust
    \underbar{cl}uster(\varname)
    \underbar{abs}orb({\it terms\/})
    d({\it newvar\/})
    \underbar{tol}erance(\num)
    maxiter(\num)
    fast
    {\it poisson\_options}}
\end{stsyntax}

The data must be \stcmd{stset}; see \stref{stset}. \stcmd{fweight}s,
\stcmd{iweight}s, and \stcmd{pweight}s are allowed on all paths;
\stcmd{pweight}s imply robust standard errors, and weighted coefficients
and standard errors agree with the corresponding weighted
\stcmd{poisson} model to at least \(10^{-6}\). \stcmd{fast} and
\stcmd{absorb()} may be combined
(section~\ref{sec:combo}). The varlist may include
factor variables and interactions on every path. Variables named in
\stcmd{absorb()} and, with \stcmd{fast}, the cluster variable must be
numeric.

\subsection{Options}\label{options}

\hangpara
\stcmd{irr} reports exponentiated coefficients (incidence-rate ratios,
equal to hazard ratios in this model).  Display only.

\hangpara
\stcmd{nolog} suppresses iteration logs.

\hangpara
\stcmd{level(\num)} sets the confidence level; the default is
\stcmd{level(95)}.

\hangpara
\stcmd{noconstant} suppresses the constant term.

\hangpara
\sloppy
\stcmd{vce({\it vcetype\/})}, \stcmd{robust}, and
\stcmd{cluster(\varname)} specify the variance estimator.
\stcmd{vce(oim)} (the default), \stcmd{vce(robust)}, and
\stcmd{vce(cluster {\it clustvar\/})} are supported on all paths
and reproduce the corresponding full \stcmd{poisson} model to at least
$10^{-6}$; the standard path additionally accepts anything
\stcmd{poisson} accepts.

\hangpara
\stcmd{fast} selects cell-accelerated exact estimation; see
section~\ref{sec:fast}.  With \stcmd{fast}, terms built only from
factor-variable pieces define the cells and enter at the cell level; at
least one such term is required.

\hangpara
\stcmd{absorb({\it terms\/})} absorbs the fixed effects of one or more
terms, each a variable name or a {\it var1\/}\stcmd{\#}{\it var2\/}
interaction; see section~\ref{sec:hdfe}.  May be combined with
\stcmd{fast}.

\hangpara
\stcmd{d({\it newvar\/})} stores the total fixed-effect contribution
of \stcmd{absorb()} in {\it newvar\/}; the default is
\stcmd{\_stpois\_fe}, replaced on each run.  The variable is used by
\stcmd{predict}.

\hangpara
\stcmd{tolerance(\num)} and \stcmd{maxiter(\num)} control the iteration
convergence criterion and cap for \stcmd{absorb()} and \stcmd{fast};
defaults are \stcmd{1e-8} and \stcmd{100}.

\medskip
\noindent
Any remaining options on the standard path are passed to
\stcmd{poisson}.

\subsection{Stored results}\label{stored-results}

\noindent\stcmd{stpois} stores the following in \stcmd{e()}:

\medskip
\noindent\textit{Scalars}\nopagebreak
\begin{tabular}{@{\hspace{1em}}ll@{}}
\stcmd{e(N)}         & number of observations (sum of the weights with \stcmd{fweight}s)\\
\stcmd{e(N\_sub)}    & number of subjects\\
\stcmd{e(N\_fail)}   & number of failures (a weighted sum with \stcmd{fweight}s)\\
\stcmd{e(risk)}      & total time at risk (a weighted sum with \stcmd{fweight}s)\\
\stcmd{e(ll)}        & log likelihood (individual-level, comparable across paths)\\
\stcmd{e(ll\_0)}     & log likelihood, constant-only model (paths without \stcmd{absorb()})\\
\stcmd{e(chi2)}      & model test statistic\\
\stcmd{e(df\_m)}     & model degrees of freedom\\
\stcmd{e(rank)}      & rank of \stcmd{e(V)}\\
\stcmd{e(k)}         & number of parameters (standard path)\\
\stcmd{e(ic)}        & number of iterations\\
\stcmd{e(converged)} & \stcmd{1} if converged, \stcmd{0} otherwise\\
\stcmd{e(N\_cells)}  & number of cells (\stcmd{fast})\\
\stcmd{e(N\_clust)}  & number of clusters (with \stcmd{vce(cluster)})\\
\end{tabular}

\medskip
\noindent\textit{Macros}\nopagebreak
\begin{tabular}{@{\hspace{1em}}ll@{}}
\stcmd{e(cmd)}       & \stcmd{stpois}\\
\stcmd{e(cmdline)}   & command as typed\\
\stcmd{e(depvar)}    & \stcmd{\_d} (the failure indicator)\\
\stcmd{e(t0)}        & \stcmd{\_t0} (the entry-time variable)\\
\stcmd{e(title)}     & title in estimation output\\
\stcmd{e(chi2type)}  & \stcmd{LR} or \stcmd{Wald}; type of model test\\
\stcmd{e(vce)}       & \stcmd{vcetype} specified in \stcmd{vce()}\\
\stcmd{e(predict)}   & \stcmd{stpois\_p} (program used to implement \stcmd{predict})\\
\stcmd{e(properties)}& \stcmd{b V}\\
\stcmd{e(fast)}      & \stcmd{fast}, if the fast path was used\\
\stcmd{e(cont\_vars)}& individual-level (continuous-involving) terms (\stcmd{fast})\\
\stcmd{e(cat\_vars)} & cell-defining categorical terms (\stcmd{fast})\\
\stcmd{e(absorb)}    & absorbed terms (\stcmd{absorb()})\\
\stcmd{e(fes\_var)}  & name of the stored fixed-effect variable (\stcmd{absorb()})\\
\stcmd{e(clustvar)}  & name of the cluster variable (with \stcmd{vce(cluster)})\\
\end{tabular}

\medskip
\noindent\textit{Matrices}\nopagebreak
\begin{tabular}{@{\hspace{1em}}ll@{}}
\stcmd{e(b)}         & coefficient vector\\
\stcmd{e(V)}         & variance--covariance matrix of the estimators\\
\end{tabular}

\medskip
\noindent\textit{Functions}\nopagebreak
\begin{tabular}{@{\hspace{1em}}ll@{}}
\stcmd{e(sample)}    & marks the estimation sample\\
\end{tabular}

All quantities are computed on the final estimation sample (after any
separation-induced drops on the \stcmd{absorb()} path), so they agree
with \stcmd{e(N)} and \stcmd{e(sample)}. Paths without absorbed effects
store a likelihood-ratio model test, with \stcmd{e(chi2type)} set to
\stcmd{LR}; with \stcmd{absorb()} the Wald test of
section~\ref{sec:hdfe} is stored, with
\stcmd{e(chi2type)} set to \stcmd{Wald}. Because every path reports the
individual-level log likelihood, likelihood-based statistics are
comparable across paths.

\begin{sloppypar}
Because \stcmd{e(predict)} is set to \stcmd{stpois\_p}, \stcmd{predict}
offers the linear predictor (\stcmd{xb}, the default, including the
log-exposure offset unless \stcmd{nooffset}), predicted numbers of
events (\stcmd{n}), hazard rates (\stcmd{hazard} or \stcmd{hr}), and
survival probabilities (\stcmd{survival}).  After \stcmd{absorb()},
\stcmd{predict} adds the stored fixed-effect contribution to the linear
predictor, so all four statistics include the absorbed effects and
match the dummy-variable \stcmd{poisson} predictions.  Postestimation
commands that operate on the coefficient vector and its variance
matrix---\stcmd{lincom}, \stcmd{test}, \stcmd{nlcom}, and
\stcmd{estimates}---work as usual on every path.  \stcmd{margins} is
available for quantities that are functions of the reported
coefficients; on the standard and \stcmd{fast} paths it and
\stcmd{predict} behave exactly as after \stcmd{poisson}, since no stored
fixed-effect variable is involved.  Note, however, that after
\stcmd{absorb()} the absorbed fixed-effect contribution enters
\stcmd{predict} through the stored \stcmd{\_stpois\_fe} variable, which
\stcmd{margins} holds fixed at its estimated values.  Margins of the
response therefore condition on the absorbed effects and their
delta-method standard errors do not propagate the sampling uncertainty
in those effects, exactly as for any fixed-effect contribution
recovered post-absorption in \stcmd{ppmlhdfe} or \stcmd{reghdfe}.
\end{sloppypar}

On the \stcmd{fast} and \stcmd{absorb()} paths, \stcmd{stpois} stores
the convergence flag \stcmd{e(converged)} and the iteration count
\stcmd{e(ic)}. If the cap \stcmd{maxiter()} is reached before the
tolerance in \stcmd{tolerance()} is met, the command still posts the
last iterate, sets \stcmd{e(converged)} to \(0\), and issues a warning,
so a noninteractive pipeline can detect the failure by inspecting
\stcmd{e(converged)}. Cell-level and non-absorbed terms that are
collinear---including collinearity with the absorbed effects---are
dropped through a generalized inverse and reflected in \stcmd{e(rank)}.
Separation is detected only for the absorbed terms, whose all-zero-event
groups are dropped before estimation
(section~\ref{sec:hdfe}); separation induced by the
non-absorbed varlist is not given a targeted drop and instead surfaces
as nonconvergence with \stcmd{e(converged)} equal to \(0\). In practice,
if \stcmd{stpois} does not converge, the first diagnostic steps are to
check \stcmd{e(ic)} against \stcmd{maxiter()}, to inspect whether any
covariate perfectly predicts the outcome within cells, and---if using
\stcmd{absorb()}---to verify that every absorbed group contains at least
one event; \stcmd{ppmlhdfe} \citep{correia2020} provides a more
comprehensive separation-detection algorithm for the HDFE case.

\section{Examples}

\label{sec:examples}

\subsection{Standard path}\label{standard-path}

The default path reproduces \stcmd{streg, distribution(exponential)}
exactly (the log likelihoods differ by a constant because \stcmd{streg}
includes the \(d_i \log t_i\) terms). We load the Stanford heart
transplant data and \stcmd{stset} them:

\begin{stlog}
. webuse stan3
(Heart transplant data)
{\smallskip}
. stset t1, failure(died) id(id)
{\smallskip}
\oom
\end{stlog}

\begin{stlog}
. stpois age posttran, nolog
{\smallskip}
Poisson event-history regression
{\smallskip}
No. of subjects =          103                  Number of obs   =      172
No. of failures =           75
Time at risk    =        31938                  LR chi2(2)      =    36.75
Log likelihood  =   -233.83404                  Prob > chi2     =   0.0000
{\smallskip}
\HLI{13}{\TOPT}\HLI{64}
          _d {\VBAR} Coefficient  Std. err.      z    P>|z|     [95\% conf. interval]
\HLI{13}{\PLUS}\HLI{64}
         age {\VBAR}   .0605817    .015235     3.98   0.000     .0307216    .0904417
    posttran {\VBAR}  -1.285206   .2377048    -5.41   0.000    -1.751099   -.8193135
       _cons {\VBAR}    -7.8376   .7094787   -11.05   0.000    -9.228153   -6.447047
\HLI{13}{\BOTT}\HLI{64}
\nullskip
\end{stlog}

\subsection{Absorbing fixed effects}\label{absorbing-fixed-effects}

\stcmd{absorb()} returns the same coefficients and standard errors as
the explicit dummy-variable model, without estimating the absorbed
effects:

\begin{stlog}
\input{ex_absorb.log.tex}\nullskip
\end{stlog}

\noindent For comparison,
\stcmd{poisson \_d age posttran i.surgery, exposure(exp)} on the same
sample gives \stcmd{age} \(=\) .0610952 (SE .014454) and
\stcmd{posttran} \(=-1.141144\) (SE .2402034): identical to all printed
digits. The Wald statistic equals the joint test of \stcmd{age} and
\stcmd{posttran} on the dummy-variable model.

\subsection{Fast estimation on a large
dataset}\label{fast-estimation-on-a-large-dataset}

The following simulation generates 500,000 subjects in 200 risk cells (4
education \(\times\) 5 region \(\times\) 10 period), with two continuous
covariates (\stcmd{age}, \stcmd{income}):

\begin{stlog}
. clear
. set obs 500000
. set seed 20260713
. gen edu    = ceil(4 * runiform())
. gen region = ceil(5 * runiform())
. gen period = ceil(10 * runiform())
. gen age    = 30 + 20 * runiform()
. gen income = rnormal(10, 2)
. gen lambda = exp(-6 + 0.03*age + 0.05*income
>     + 0.3*(edu==2) + 0.5*(edu==3) - 0.2*(edu==4)
>     + 0.1*(region==2) - 0.2*(region==3) + 0.2*(period >= 6))
. gen ptime  = 0.5 + 0.5 * runiform()
. gen failed = (runiform() < 1 - exp(-lambda * ptime))
. gen id2    = _n
. stset ptime, failure(failed) id(id2)
{\smallskip}
\oom
\end{stlog}

\begin{stlog}
. stpois age income i.edu i.region i.period, fast nolog
  (fast) 2 individual-level + 16 cell-level parameters...
  (fast) 500000 obs, 200 cells, 5 Newton iterations, converged=1
{\smallskip}
Poisson EHA {\EMDASH} fast (cell-accelerated exact MLE)
{\smallskip}
No. of subjects =      500,000                  Number of obs   =  500,000
No. of failures =        6,759
Time at risk    =       375045                  LR chi2(18)     =   850.69
Log likelihood  = -35297.23308                  Prob > chi2     =   0.0000
{\smallskip}
\HLI{13}{\TOPT}\HLI{64}
          _d {\VBAR} Coefficient  Std. err.      z    P>|z|     [95\% conf. interval]
\HLI{13}{\PLUS}\HLI{64}
         age {\VBAR}   .0265379    .002121    12.51   0.000     .0223808    .0306951
      income {\VBAR}   .0518263   .0060872     8.51   0.000     .0398956    .0637571
             {\VBAR}
         edu {\VBAR}
          2  {\VBAR}   .3070999   .0353078     8.70   0.000      .237898    .3763019
          3  {\VBAR}   .5094171   .0339083    15.02   0.000      .442958    .5758762
          4  {\VBAR}  -.1907793   .0398372    -4.79   0.000    -.2688587   -.1126998
             {\VBAR}
\oom
       _cons {\VBAR}  -5.885082   .1191483   -49.39   0.000    -6.118608   -5.651556
\HLI{13}{\BOTT}\HLI{64}
\nullskip
\end{stlog}

\noindent Running the same model on the standard path returns identical
output---the same coefficients, standard errors, and log likelihood to
all digits---in five Newton iterations versus a full \stcmd{poisson}
maximization; the timing comparison is in
section~\ref{sec:validation}.

\subsection{Interactions}\label{interactions}

All interaction types are accepted with \stcmd{fast}; terms involving a
continuous piece, such as the \stcmd{edu\#c.age} terms produced by
\stcmd{c.age\#\#i.edu} below, are estimated at the individual level,
while the categorical main effects and any categorical\#categorical
interactions run at the cell level:

\begin{stlog}
. stpois c.age\#\#i.edu income i.region i.period, fast nolog
  (fast) 5 individual-level + 16 cell-level parameters...
  (fast) 500000 obs, 200 cells, 5 Newton iterations, converged=1
{\smallskip}
Poisson EHA {\EMDASH} fast (cell-accelerated exact MLE)
{\smallskip}
No. of subjects =      500,000                  Number of obs   =  500,000
No. of failures =        6,759
Time at risk    =       375045                  LR chi2(21)     =   851.30
Log likelihood  = -35296.92816                  Prob > chi2     =   0.0000
{\smallskip}
\HLI{13}{\TOPT}\HLI{64}
          _d {\VBAR} Coefficient  Std. err.      z    P>|z|     [95\% conf. interval]
\HLI{13}{\PLUS}\HLI{64}
         age {\VBAR}   .0255056   .0046717     5.46   0.000     .0163492    .0346621
             {\VBAR}
   edu\#c.age {\VBAR}
          2  {\VBAR}  -.0009954   .0061531    -0.16   0.871    -.0130552    .0110645
          3  {\VBAR}   .0021302   .0059106     0.36   0.719    -.0094545    .0137148
          4  {\VBAR}   .0034035   .0069472     0.49   0.624    -.0102128    .0170198
             {\VBAR}
\oom
       _cons {\VBAR}  -5.842836   .2075639   -28.15   0.000    -6.249653   -5.436018
\HLI{13}{\BOTT}\HLI{64}
\nullskip
\end{stlog}

\noindent Interacted fixed effects are absorbed with
\stcmd{absorb(edu\#region)}, which matches \stcmd{poisson} with explicit
\stcmd{i.edu\#i.region} dummies on coefficients and standard errors.

\subsection{Combining fast with absorb(), and
predictions}\label{combining-fast-with-absorb-and-predictions}

Returning to the Stanford heart-transplant data, the two acceleration
mechanisms combine, and both absorption engines store the fixed-effect
contribution that \stcmd{predict} uses:

\begin{stlog}
. webuse stan3, clear
(Heart transplant data)
{\smallskip}
. stset t1, failure(died) id(id)
{\smallskip}
\oom
\end{stlog}

\begin{stlog}
. stpois age i.posttran, fast absorb(surgery) nolog
  (fast) 1 individual-level + 1 cell-level parameters...
  (fast) 172 obs, 2 cells, 5 Newton iterations, converged=1
{\smallskip}
Poisson EHA {\EMDASH} fast + HDFE (absorb: surgery)
{\smallskip}
No. of subjects =          103                  Number of obs   =      172
No. of failures =           75
Time at risk    =        31938                  Wald chi2(2)    =    35.17
Log likelihood  =   -227.61949                  Prob > chi2     =   0.0000
{\smallskip}
\HLI{13}{\TOPT}\HLI{64}
          _d {\VBAR} Coefficient  Std. err.      z    P>|z|     [95\% conf. interval]
\HLI{13}{\PLUS}\HLI{64}
         age {\VBAR}   .0610952    .014454     4.23   0.000     .0327659    .0894244
  1.posttran {\VBAR}  -1.141144   .2402034    -4.75   0.000    -1.611934   -.6703538
\HLI{13}{\BOTT}\HLI{64}
\nullskip
\end{stlog}

\noindent The same coefficients and standard errors are returned by the
IRLS engine (\stcmd{absorb()} without \stcmd{fast}) and by
\stcmd{poisson} with explicit \stcmd{i.surgery} dummies. Predictions
after \stcmd{absorb()} include the absorbed effects:

\begin{stlog}
. stpois age posttran, absorb(surgery) nolog
  (HDFE: 6 IRLS iterations, converged=1)
{\smallskip}
Poisson EHA {\EMDASH} HDFE (absorb: surgery)
{\smallskip}
No. of subjects =          103                  Number of obs   =      172
No. of failures =           75
Time at risk    =        31938                  Wald chi2(2)    =    35.17
Log likelihood  =   -227.61949                  Prob > chi2     =   0.0000
{\smallskip}
\HLI{13}{\TOPT}\HLI{64}
          _d {\VBAR} Coefficient  Std. err.      z    P>|z|     [95\% conf. interval]
\HLI{13}{\PLUS}\HLI{64}
         age {\VBAR}   .0610952    .014454     4.23   0.000     .0327659    .0894244
    posttran {\VBAR}  -1.141144   .2402034    -4.75   0.000    -1.611934   -.6703538
\HLI{13}{\BOTT}\HLI{64}
{\smallskip}
. predict double haz, hazard
{\smallskip}
. summarize haz
{\smallskip}
    Variable {\VBAR}        Obs        Mean    Std. dev.       Min        Max
\HLI{13}{\PLUS}\HLI{57}
         haz {\VBAR}        172    .0050208    .0039522   .0003224   .0212179
{\smallskip}
\nullskip
\end{stlog}

\subsection{Episode-split workflow}\label{episode-split-workflow}

In applications, the typical sequence splits spells on calendar time,
age, or duration and then estimates piecewise-constant hazard models:

\begin{stlog}
. stsplit fu, at(1(1)30)
. stpois i.fu age income i.edu, irr
. stpois i.fu age income i.edu, fast
. stpois age income, absorb(fu edu region) vce(cluster region)
\end{stlog}

\noindent The last command absorbs three fixed-effect terms
simultaneously and clusters the standard errors; the full output for
this workflow is generated by the example script that accompanies the
package.

\section{Relation to existing commands}

\label{sec:related}

Table~\ref{tab:features} places \stcmd{stpois} among the
Stata commands that fit closely related models. For one-way fixed
effects, the official \stcmd{xtpoisson, fe} conditional estimator is
dependency-free and consistent, but it conditions the fixed effects out
of the likelihood rather than absorbing them, so it does not extend to
multi-way fixed effects, does not recover or predict the fixed-effect
contribution, and its clustered-variance conventions differ from those
of \stcmd{poisson} with explicit dummies; \stcmd{stpois, absorb()}
targets exactly the cases---multi-way absorption, recovered effects, and
\stcmd{poisson}-matching robust and clustered standard errors---in which
the conditional estimator does not apply. Against \stcmd{ppmlhdfe}, the
closest command for the absorption path, \stcmd{stpois} adds
\stcmd{stset} integration and the cell-accelerated fast path while
requiring only Stata~14 or later and carrying no external
dependencies, at the cost of a narrower separation-detection scheme
(section~\ref{sec:hdfe}).

\begin{table}[h!]
\caption{Feature coverage relative to related commands ($\bullet$~supported
natively, $\circ$~available only indirectly, e.g.\ through explicit
dummies or with restrictions, blank~not available)}
\label{tab:features}
\begin{center}
{\footnotesize
\setlength{\tabcolsep}{4pt}
\begin{tabular}{lcccccc}
\hline
\noalign{\smallskip}
Feature & \stcmd{poisson} & \stcmd{streg} & \stcmd{stcox}
  & \parbox{1.3cm}{\centering\stcmd{xtpois-}\\\stcmd{son, fe}}
  & \stcmd{ppmlhdfe} & \stcmd{stpois} \\
\noalign{\smallskip}
\hline
\noalign{\smallskip}
\stcmd{stset}-native interface   &        &$\bullet$&$\bullet$&        &        &$\bullet$\\
Survival headers/predict         &        &$\bullet$&$\bullet$&        &        &$\bullet$\\
One-way fixed effects            &$\circ$ &        &        &$\bullet$&$\bullet$&$\bullet$\\
Multi-way fixed effects          &$\circ$ &        &        &        &$\bullet$&$\bullet$\\
Cell-accelerated continuous fit  &        &        &        &        &        &$\bullet$\\
Robust/cluster VCE               &$\bullet$&$\bullet$&$\bullet$&$\bullet$&$\bullet$&$\bullet$\\
fw/iw/pw weights                 &$\bullet$&$\bullet$&$\bullet$&$\circ$ &$\bullet$&$\bullet$\\
Separation detection             &        &        &        &        &$\bullet$&$\circ$ \\
Recover/predict FE contribution  &        &        &        &        &$\bullet$&$\bullet$\\
No external dependencies         &$\bullet$&$\bullet$&$\bullet$&$\bullet$&        &$\bullet$\\
\noalign{\smallskip}
\hline
\end{tabular}
}
\end{center}
\end{table}

\section{Validation and performance}

\label{sec:validation}

\subsection{Numerical agreement with streg and
poisson}\label{numerical-agreement-with-streg-and-poisson}

The standard and fast estimators are algebraically exact in that they
solve the same likelihood as \stcmd{poisson} with no mathematical
approximation; the \stcmd{absorb()} estimator converges to the same
maximum-likelihood optimum to within its specified tolerance.
``Exact'\,' throughout this article therefore means agreement with the
\stcmd{streg}/\stcmd{poisson} benchmarks to at least \(10^{-6}\) on
coefficients \emph{and} all three variance estimators, with observed
agreement typically on the order of \(10^{-8}\); it does not assert
bit-for-bit equality of two independent maximizers. The package ships
with a test suite (\stcmd{test\_stpois.do} and
\stcmd{test\_fast\_exact.do} in the \stcmd{tests} directory) that runs
on the \stcmd{stan3} heart-transplant data and on simulated data---not
on register data, which cannot leave the secure environments that
motivate the command. Each exactness claim is held to a numerical
benchmark:

\begin{enumerate}
\item The \textbf{standard path} matches \stcmd{streg, distribution(exponential)} on every coefficient and standard error to $10^{-6}$.  On the \stcmd{stan3} data, both report $\hat{\beta}_{\rm age} = 0.0605817$ (SE $0.0152350$) and $\hat{\beta}_{\rm posttran} = -1.2852060$ (SE $0.2377048$); the log likelihoods differ only by the fixed $\sum_i d_i \log t_i$ term that \stcmd{streg} includes and the Poisson formulation drops.
\item \sloppy \textbf{absorb()} matches the explicit dummy-variable \stcmd{poisson} model on coefficients \emph{and} standard errors to $10^{-6}$ (observed agreement is typically on the order of $10^{-8}$), for the observed-information, robust, and cluster--robust variance estimators.  The check covers single fixed effects, interacted fixed effects (\stcmd{edu\#region} against explicit \stcmd{i.edu\#i.region} dummies), multiple simultaneous fixed-effect terms (\stcmd{absorb(edu region period)}), and factor-variable interactions in the varlist.  The Wald model statistic matches the joint test on the dummy-variable model, and a separation test verifies that the header statistics are computed on the post-drop estimation sample.
\item \textbf{fast} matches the full \stcmd{poisson} model on coefficients, standard errors (observed-information, robust, and clustered), and log likelihood to $10^{-6}$.  The check covers continuous\#continuous, continuous\#categorical, and categorical\#categorical interactions, all-categorical models, and a two-cell setting in which the continuous covariate is constructed to have exactly zero between-cell variation, so that its coefficient is identified purely from within-cell variation.
\item \textbf{Weights} (frequency, importance, and sampling) on the fast and absorbed paths match the corresponding weighted \stcmd{poisson} model on coefficients and standard errors to $10^{-6}$, with \stcmd{e(N)} following Stata's conventions and \stcmd{pweight}s implying robust standard errors.
\item \textbf{fast + absorb()} matches \stcmd{poisson} with explicit dummies for both the cell-level terms and the absorbed effects---coefficients and observed-information, robust, and clustered standard errors to $10^{-6}$---including multi-term absorbs and weighted estimation.
\item \textbf{predict} after \stcmd{absorb()} reproduces the dummy-variable \stcmd{poisson} predictions of expected events, and the hazard equals the exponentiated no-offset linear predictor, on both the IRLS and the fast+absorb engines.
\end{enumerate}

All tests pass on Stata~19 (BE); the command requires
version~14. The suite was run on a clean Stata
installation with no other user-contributed packages present, confirming
that \stcmd{stpois} draws on no external dependencies.

\subsection{Agreement with ppmlhdfe and with
collapse+poisson}\label{agreement-with-ppmlhdfe-and-with-collapsepoisson}

\label{sec:headtohead}

Because \stcmd{absorb()} targets the same problem as \stcmd{ppmlhdfe}
and \stcmd{fast} generalizes the classical collapse-and-fit workflow,
the accompanying script \stcmd{benchmark\_ppmlhdfe.do} checks
\stcmd{stpois} directly against both on a \(10^6\)-record design with a
1,000-level fixed effect. The point estimates are identical: on the
\stcmd{absorb()} path \stcmd{stpois} and \stcmd{ppmlhdfe} agree on
\(\hat{\beta}_{\rm age}\) to \(2\times10^{-13}\), and on the combined
\stcmd{fast}\({}+{}\)\stcmd{absorb()} path to \(1\times10^{-14}\).
Compared like for like, the standard errors agree to about
\(3\times10^{-9}\): \stcmd{ppmlhdfe} reports robust standard errors by
default, as is standard for Poisson pseudo-maximum-likelihood, and these
match \stcmd{stpois} under \stcmd{vce(robust)} and \stcmd{vce(cluster)}
to that tolerance. \stcmd{stpois} additionally provides \stcmd{vce(oim)}
as its default, reproducing \stcmd{poisson}'s classical
observed-information standard errors---an option \stcmd{ppmlhdfe} does
not offer (as of version 2.4.1, the version benchmarked here). Against
the folk baseline of collapsing to cells and fitting weighted
\stcmd{poisson} on the sums, the \stcmd{fast} path reproduces the
all-categorical estimates and their standard errors to machine precision
(\(<10^{-15}\))---as it must, since for this case (no absorbed effects,
default VCE) \stcmd{fast} \emph{is} that collapse-and-fit workflow,
performed internally so the user neither writes the \stcmd{collapse} nor
leaves the \stcmd{stset} data (section~\ref{sec:fast}).
What \stcmd{fast} adds beyond the manual workflow is support for
continuous covariates, for which no sufficient cell reduction exists
(section~\ref{sec:suff}): the command accelerates
through the cell structure while retaining the continuous terms in the
individual-level tilted moments. On timing (medians of ten runs, same
machine as the tables), \stcmd{stpois, absorb()} fits the 1,000-level
design of Table~\ref{tab:hdfe} in 2.9~ s
against \stcmd{ppmlhdfe}'s 4.6~ s, and the combined
\stcmd{fast}\({}+{}\)\stcmd{absorb()} path in 5.0~ s
against 11.6~ s; \stcmd{stpois} is thus competitive with,
and here faster than, the established command while adding the
\stcmd{stset} interface and survival postestimation. On the
all-categorical design, the internal collapse route fits in
0.5~ s versus 0.1~ s for hand-collapsing and
calling \stcmd{poisson}---roughly 5\(\times\) slower---the gap arising
from \stcmd{stpois}'s fixed overhead of sample marking, term parsing,
and result posting; what the command adds is the \stcmd{stset} workflow
and continuous-covariate support without requiring the user to write a
\stcmd{collapse}.

\subsection{Timing benchmarks}\label{timing-benchmarks}

All timings are medians over ten runs on Stata~19 (BE,
Apple silicon, warm cache), produced by the benchmark scripts in the
package repository, which write their results to comma-separated files
versioned alongside the paper; the tables below reproduce those files.
Reported speedup ratios are computed from the unrounded median timings,
so they need not equal the ratio of the rounded cells shown. In every
comparison the estimates agree across methods to the tolerances
above---the scripts assert this inside the timing loops---so the tables
compare identical output.

Table~\ref{tab:mixed} times a mixed model with two
continuous and 16 categorical parameters (education, region, period),
comparing \stcmd{streg, distribution(exponential)}, the \stcmd{stpois}
standard path (a \stcmd{poisson} call), and \stcmd{stpois, fast}.

\begin{table}[h!]
\caption{Mixed continuous $+$ categorical model, seconds per fit.
DGP: 200 cells (4 education $\times$ 5 region $\times$ 10 period),
two continuous covariates (age, income), $\approx 1\%$ event rate.
Medians of 10 runs, Stata~19 BE, Apple silicon, warm cache.}
\label{tab:mixed}
\begin{center}
\begin{tabular}{rrrrr}
\hline
\noalign{\smallskip}
$n$ & \stcmd{streg} & \stcmd{stpois} (standard) & \stcmd{fast}
  & speedup vs.\ \stcmd{streg} \\
\noalign{\smallskip}
\hline
\noalign{\smallskip}
100,000 & 0.37 & 0.34 & 0.11 & 3.3$\times$ \\
1,000,000 & 2.73 & 2.52 & 1.10 & 2.5$\times$ \\
5,000,000 & 13.26 & 12.41 & 6.72 & 2.0$\times$ \\
\noalign{\smallskip}
\hline
\end{tabular}
\end{center}
\end{table}

Table~\ref{tab:split} times the canonical episode-split
workflow: 100,000 subjects observed for up to 30 years are split into
one-year episodes with \stcmd{stsplit} (1,888,654 episode records), and
a piecewise-constant hazard model with 30 follow-up bands, education,
region, and two continuous covariates is estimated on the episode file.

\begin{table}[h!]
\caption{Episode-split workflow (1,888,654 episode records).
DGP: 100,000 subjects split into annual episodes over up to 30 years;
piecewise-constant hazard with 30 follow-up bands, education, region,
and two continuous covariates.
Medians of 10 runs, Stata~19 BE, Apple silicon, warm cache.}
\label{tab:split}
\begin{center}
\begin{tabular}{lrr}
\hline
\noalign{\smallskip}
Method & Time & speedup vs.\ \stcmd{streg} \\
\noalign{\smallskip}
\hline
\noalign{\smallskip}
\stcmd{streg, distribution(exponential)} & 5.84 s & --- \\
\stcmd{stpois} (standard)                & 5.12 s & 1.1$\times$ \\
\stcmd{stpois, fast}                     & 2.19 s   & 2.7$\times$ \\
\noalign{\smallskip}
\hline
\end{tabular}
\end{center}
\end{table}

Table~\ref{tab:allcat} times an all-categorical rate
model (education \(\times\) region \(\times\) period, 200 cells), the
case where the sufficiency result of
section~\ref{sec:suff} lets \stcmd{fast} fit on the
collapsed cell sums alone through the internal
\stcmd{collapse}\({}+{}\)\stcmd{poisson} route.

\begin{table}[h!]
\caption{All-categorical rate model, seconds per fit.
DGP: 200 cells (4 education $\times$ 5 region $\times$ 10 period),
no continuous covariates, $\approx 1\%$ event rate.
Medians of 10 runs, Stata~19 BE, Apple silicon, warm cache.
For this case \stcmd{fast} internally routes through \stcmd{collapse}$+$\stcmd{poisson}
on the 200 cell sums; hand-collapsing and calling \stcmd{poisson} directly
is $\approx 5\times$ faster (section~\ref{sec:headtohead}).}
\label{tab:allcat}
\begin{center}
\begin{tabular}{rrrrr}
\hline
\noalign{\smallskip}
$n$ & \stcmd{streg} & \stcmd{stpois} (standard) & \stcmd{fast}
  & speedup vs.\ \stcmd{streg} \\
\noalign{\smallskip}
\hline
\noalign{\smallskip}
1,000,000 & 3.60 & 3.43 & 0.60 & 6.0$\times$ \\
5,000,000 & 13.05 & 12.06 & 2.86 & 4.6$\times$ \\
\noalign{\smallskip}
\hline
\end{tabular}
\end{center}
\end{table}

In tables~\ref{tab:mixed}--\ref{tab:allcat} the
complexity analysis implies that the remaining \stcmd{fast} cost is
dominated by the single aggregating pass over the microdata: the
per-iteration work is \(O(np^2)\) in the few individual-level columns
only (Table~\ref{tab:mixed}), and for the
all-categorical model (Table~\ref{tab:allcat}) it
reduces to one \stcmd{collapse} plus a \stcmd{poisson} fit on a few
hundred cells. The gains are therefore largest relative to estimators
that must pass over the microdata repeatedly, which is what the
fixed-effects benchmark isolates. Table~\ref{tab:hdfe}
estimates a model with two continuous covariates and a single fixed
effect with 1,000 levels on one million observations, comparing
\stcmd{poisson} with explicit dummies, \stcmd{stpois, absorb()}, and
\stcmd{ppmlhdfe} on identical output.

\begin{table}[h!]
\caption{High-dimensional fixed effect: 1,000 levels, $n = 10^6$,
two continuous covariates.
\stcmd{ppmlhdfe} reports robust standard errors by default and offers no
observed-information option (section~\ref{sec:headtohead}).
Rows 1--3: medians of 10 runs from \texttt{benchmark\_hdfe.do};
row 4: medians of 10 runs from \texttt{benchmark\_ppmlhdfe.do}
on the same DGP (robust SEs; \stcmd{stpois} SE omitted for brevity).
Stata~19 BE, Apple silicon, warm cache.}
\label{tab:hdfe}
\begin{center}
\begin{tabular}{lrrr}
\hline
\noalign{\smallskip}
Method & Time & $\hat{\beta}_{\rm age}$ & SE \\
\noalign{\smallskip}
\hline
\noalign{\smallskip}
\stcmd{poisson} with 1,000 dummies       & 34.31 s & 0.0266067 & 0.0017000 \\
\stcmd{stpois, absorb(g)}                & 2.90 s & 0.0266067 & 0.0017000 \\
\stcmd{ppmlhdfe, absorb(g)}              & 4.59 s  & 0.0266067 & --- \\
\stcmd{stpois, fast absorb(g)}           & 5.03 s   & \multicolumn{2}{c}{(see section~\ref{sec:headtohead})} \\
\stcmd{ppmlhdfe, absorb(g)} [fast comp.] & 11.64 s & \multicolumn{2}{c}{} \\
\noalign{\smallskip}
\hline
\end{tabular}
\end{center}
\end{table}

Nearly a 12\(\times\) saving at identical output; the dummy-variable
model's cost grows roughly quadratically in the number of levels while
the absorption cost grows only linearly, so the complexity analysis
implies that the gap should widen further in higher-dimensional
settings. For register applications with \(10^7\) or more episodes,
detailed period \(\times\) education \(\times\) region structures, and
high-dimensional nuisance effects---a regime beyond the scale of the
benchmarks reported here---cell acceleration and fixed-effect absorption
together may reduce runtimes from minutes to seconds, which would change
what is practical to do interactively: specification searches, bootstrap
loops, or estimation inside secure environments with tight session
limits.

\section{Conclusion}

\stcmd{stpois} packages three things that belong together but have
required separate workflows: the Poisson formulation of event-history
models on \stcmd{stset} data, absorption of high-dimensional fixed
effects without external dependencies, and cell-accelerated estimation
that retains continuous covariates. The standard and fast paths are
algebraically exact---they solve the same likelihood as \stcmd{poisson}
with no mathematical approximation---and the \stcmd{absorb()} path
converges to the same maximum-likelihood optimum to within its specified
tolerance; the test suite holds coefficients \emph{and} variance
estimators, not merely point estimates, to the \stcmd{streg} and
\stcmd{poisson} benchmarks at \(10^{-6}\) (typically \(\sim 10^{-8}\)),
and factor-variable interactions are supported throughout, including
interacted fixed effects. Weights are accepted on every path, the fast
and absorbed paths can be combined, and predictions after
\stcmd{absorb()} include the stored fixed-effect contributions. The
command, its help file, the test suite, and the benchmark scripts are
available from the author's repository and are installable with
\stcmd{net install}.

\section{Acknowledgments}

TBW

\bibliographystyle{plainnat}
\bibliography{stpois}

\bigskip\noindent\textbf{About the author.} Andreas Ljungström is a PhD
student at the Swedish Institute for Social Research (SOFI), Stockholm
University. His work concerns social stratification broadly, but with a
special focus on integration in education and early labor market
outcomes using full-population Swedish register data.

\end{document}
